\section{Inequacions}

\begin{enumerate}
    \item Trobau les solucions de $|x-|x-2||<3$:\\
Per $x\ge 2$:

$|x-(x-2)|<3 \implies |2|=2<3, \forall x\in \mathbb{R}^{\ge2}: |x-|x-2||<3$

Per $x<2$: 

$|x+x-2|<3 \implies |2x-2|<3$, hi ha dues opcións:\\

$2x-2\ge0 \implies 2x-2=3 \implies x=\frac{5}{2}$
$2x-2<0 \implies -2x+2=3 \implies x=-\frac{1}{2}$

Ara agafant $x\in (-1/2, 2)$ (per exemple $x=0$) i a $x=2$, la inequació es cumpleix, mentre que a $(-\infty, -1/2 )$ i $(2, \infty)$ no es cumpleix. $x\in (-1/2, 2]\\$
Ara feim la unió dels dos conjunts: $x\in (-1/2, 2]\cup [2, \infty)= (-1/2, \infty)$

    \item Conjunt de punts que satisfà $x^2-ax\le2$: 
    $x^2-ax-2=0 \implies x=\frac{a\pm \sqrt{a^2-4(-2)}}{2} \implies x_1=\frac{a- \sqrt{a^2+8}}{2}, x_2=\frac{a+ \sqrt{a^2+8}}{2}$ 
Ara, agafant el punt mitjá entre $x_1, x_2$: $(\frac{a+ \sqrt{a^2+8}}{2}+\frac{a- \sqrt{a^2+8}}{2})/2=a/2$, substituint a la funció: $a^2/4 - a^2/2-2 = -a^2/2-2=-(a^2/2+2)\le 0$. Com som a una inecuació de segon ordre només es cumpleix per $x\in [\frac{a- \sqrt{a^2+8}}{2},\frac{a+ \sqrt{a^2+8}}{2}]$
    


    
\end{enumerate}